\subsection{酉性}\label{hermcon}

最后，我们解释构造出的辛同胚算符的酉性。我们的讨论基于第 \ref{hermitian} 节的框架。要定义伴随，需要算子空间上的 Hermite 内积。这里我们可以分两步。首先，对偶或转置算符可由把双向图旋转 $\pi$（或等价地，把图从下往上读）得到。与辛同胚 $\varphi$ 相伴的界面的转置，来自 $\varphi^{-1}$。然后，若 $A$-模型有反辛对合 $\tau$，且膜 $\B_1,\B_2$ 为 $\tau$-不变，则 $\Hom(\B_1,\B_2)$ 有 Hermite 度量，伴随由
\be\label{wedj}\O^\dagger=\Theta_\tau\O^*\Theta_\tau\ee
给出，其中 $\O^*$ 是转置算子。

在我们的量子化情形，$\B_1=\B$（支撑在 $M$ 上的拉格朗日膜），$\B_2=\B_\cc$。若 $\varphi$ 是 $M$ 的辛同胚 $\varphi_r$ 的解析延拓，则 $\varphi$ 为 $\tau$-不变，即 $\tau\varphi\tau=\varphi$。这蕴含 $\tau\varphi^{-1}\tau=\varphi^{-1}$。因此 $\O^\dagger$ 来自界面 $\I_{\varphi^{-1}}$。于是 $\O^\dagger\O$ 来自 $\I_{\varphi^{-1}}\cdot\I_\varphi=\I_{\text{平凡}}$。平凡界面配上平凡端点，给出恒等算符乘以某个模 1 的复数。最后，算符 $\O^\dagger\O$ 非负，故该复数必为 1：
\be\label{kunu}\O^\dagger\O=1.\ee
因此构造出的辛同胚算符都是酉的。
